Em Santos, \textit{A Terra Dá, A Terra Quer}, p. 42: \enquote{no inverno íamos para a roça, mas durante o período sem chuva íamos para as casas de farinha e para os engenhos porque lá trabalhávamos na sombra. Passávamos dois meses na sombra, festejando. Nas casas de farinha não se ganhava dinheiro, assim como nos engenhos. Ganhava-se farinha, tapioca, massa, crueira. As pessoas que não tinham cana plantada, mas sabiam fazer alguma coisa, iam trabalhar para ganhar rapadura, mel, vinagre, açúcar ou cachaça. Quem não tinha engenho ia para receber o que precisava. Todo mundo podia se abastecer.}